% Môn: Vật lý
% Lớp: 10
% Chương: Chương I. Mở đầu
% Bài: L10C1 Bài 2. Các quy tắc an toàn trong phòng thực hành Vật lí · Xây dựng quy trình thực hành an toàn
% Số câu: 185
% App đọc file này trực tiếp — sửa rồi Commit trên GitHub, bấm Đồng bộ GitHub .tex

% L10C1 Bài 2. Các quy tắc an toàn trong phòng thực hành Vật lí



% ===== Câu 48 =====
% ID: L10C1B2-127-TN
% Mức: NB
\dangbt{Xây dựng quy trình thực hành an toàn}
\begin{ex}
% ID: L10C1B2-127-TN
% Mức: NB
Trước khi bắt đầu thí nghiệm liên quan đến quy trình thực hành an toàn, hành động nào an toàn nhất?
\choice
{\True lập trình tự chuẩn bị, thao tác, giám sát, kết thúc và xử lí sự cố}
{chỉ ghi các bước thao tác mà không có bước kiểm tra và ứng phó}
{tiếp tục thao tác rồi báo cáo sau}
{nhờ bạn khác làm thay mà không kiểm tra điều kiện an toàn}
\loigiai{
Hành động đúng là “lập trình tự chuẩn bị, thao tác, giám sát, kết thúc và xử lí sự cố” vì bảo đảm công việc có kiểm soát từ trước, trong và sau thí nghiệm. Chọn A.
}
\end{ex}

% ===== Câu 49 =====
% ID: L10C1B2-129-TN
% Mức: NB
\begin{ex}
% ID: L10C1B2-129-TN
% Mức: NB
Trong lúc thay đổi cách bố trí dụng cụ liên quan đến quy trình thực hành an toàn, hành động nào an toàn nhất?
\choice
{tiếp tục thao tác rồi báo cáo sau}
{nhờ bạn khác làm thay mà không kiểm tra điều kiện an toàn}
{\True lập trình tự chuẩn bị, thao tác, giám sát, kết thúc và xử lí sự cố}
{chỉ ghi các bước thao tác mà không có bước kiểm tra và ứng phó}
\loigiai{
Hành động đúng là “lập trình tự chuẩn bị, thao tác, giám sát, kết thúc và xử lí sự cố” vì bảo đảm công việc có kiểm soát từ trước, trong và sau thí nghiệm. Chọn C.
}
\end{ex}

% ===== Câu 54 =====
% ID: L10C1B2-130-TN
% Mức: TH
\begin{ex}
% ID: L10C1B2-130-TN
% Mức: TH
Sau khi hoàn thành phép đo liên quan đến quy trình thực hành an toàn, hành động nào an toàn nhất?
\choice
{nhờ bạn khác làm thay mà không kiểm tra điều kiện an toàn}
{\True lập trình tự chuẩn bị, thao tác, giám sát, kết thúc và xử lí sự cố}
{chỉ ghi các bước thao tác mà không có bước kiểm tra và ứng phó}
{tiếp tục thao tác rồi báo cáo sau}
\loigiai{
Hành động đúng là “lập trình tự chuẩn bị, thao tác, giám sát, kết thúc và xử lí sự cố” vì bảo đảm công việc có kiểm soát từ trước, trong và sau thí nghiệm. Chọn B.
}
\end{ex}

% ===== Câu 55 =====
% ID: L10C1B2-131-TN
% Mức: TH
\begin{ex}
% ID: L10C1B2-131-TN
% Mức: TH
Khi một bạn chưa đọc hướng dẫn liên quan đến quy trình thực hành an toàn, hành động nào an toàn nhất?
\choice
{\True lập trình tự chuẩn bị, thao tác, giám sát, kết thúc và xử lí sự cố}
{chỉ ghi các bước thao tác mà không có bước kiểm tra và ứng phó}
{tiếp tục thao tác rồi báo cáo sau}
{nhờ bạn khác làm thay mà không kiểm tra điều kiện an toàn}
\loigiai{
Hành động đúng là “lập trình tự chuẩn bị, thao tác, giám sát, kết thúc và xử lí sự cố” vì bảo đảm công việc có kiểm soát từ trước, trong và sau thí nghiệm. Chọn A.
}
\end{ex}

% ===== Câu 56 =====
% ID: L10C1B2-132-TN
% Mức: TH
\begin{ex}
% ID: L10C1B2-132-TN
% Mức: TH
Khi khu vực làm việc bị ướt liên quan đến quy trình thực hành an toàn, hành động nào an toàn nhất?
\choice
{chỉ ghi các bước thao tác mà không có bước kiểm tra và ứng phó}
{tiếp tục thao tác rồi báo cáo sau}
{nhờ bạn khác làm thay mà không kiểm tra điều kiện an toàn}
{\True lập trình tự chuẩn bị, thao tác, giám sát, kết thúc và xử lí sự cố}
\loigiai{
Hành động đúng là “lập trình tự chuẩn bị, thao tác, giám sát, kết thúc và xử lí sự cố” vì bảo đảm công việc có kiểm soát từ trước, trong và sau thí nghiệm. Chọn D.
}
\end{ex}

% ===== Câu 57 =====
% ID: L10C1B2-133-TN
% Mức: TH
\begin{ex}
% ID: L10C1B2-133-TN
% Mức: TH
Khi cần rời bàn thí nghiệm liên quan đến quy trình thực hành an toàn, hành động nào an toàn nhất?
\choice
{tiếp tục thao tác rồi báo cáo sau}
{nhờ bạn khác làm thay mà không kiểm tra điều kiện an toàn}
{\True lập trình tự chuẩn bị, thao tác, giám sát, kết thúc và xử lí sự cố}
{chỉ ghi các bước thao tác mà không có bước kiểm tra và ứng phó}
\loigiai{
Hành động đúng là “lập trình tự chuẩn bị, thao tác, giám sát, kết thúc và xử lí sự cố” vì bảo đảm công việc có kiểm soát từ trước, trong và sau thí nghiệm. Chọn C.
}
\end{ex}

% ===== Câu 58 =====
% ID: L10C1B2-135-TN
% Mức: VD
\begin{ex}
% ID: L10C1B2-135-TN
% Mức: VD
Khi kết quả đo khác dự kiến liên quan đến quy trình thực hành an toàn, hành động nào an toàn nhất?
\choice
{\True lập trình tự chuẩn bị, thao tác, giám sát, kết thúc và xử lí sự cố}
{chỉ ghi các bước thao tác mà không có bước kiểm tra và ứng phó}
{tiếp tục thao tác rồi báo cáo sau}
{nhờ bạn khác làm thay mà không kiểm tra điều kiện an toàn}
\loigiai{
Hành động đúng là “lập trình tự chuẩn bị, thao tác, giám sát, kết thúc và xử lí sự cố” vì bảo đảm công việc có kiểm soát từ trước, trong và sau thí nghiệm. Chọn A.
}
\end{ex}

% ===== Câu 59 =====
% ID: L10C1B2-136-TN
% Mức: VD
\begin{ex}
% ID: L10C1B2-136-TN
% Mức: VD
Khi giáo viên yêu cầu dừng thao tác liên quan đến quy trình thực hành an toàn, hành động nào an toàn nhất?
\choice
{chỉ ghi các bước thao tác mà không có bước kiểm tra và ứng phó}
{tiếp tục thao tác rồi báo cáo sau}
{nhờ bạn khác làm thay mà không kiểm tra điều kiện an toàn}
{\True lập trình tự chuẩn bị, thao tác, giám sát, kết thúc và xử lí sự cố}
\loigiai{
Hành động đúng là “lập trình tự chuẩn bị, thao tác, giám sát, kết thúc và xử lí sự cố” vì bảo đảm công việc có kiểm soát từ trước, trong và sau thí nghiệm. Chọn D.
}
\end{ex}

% ===== Câu 120 =====
% ID: L10C1B2-37-TL
% Mức: TH
\begin{bt}
% ID: L10C1B2-37-TL
% Mức: TH
Trình bày nguyên tắc cốt lõi khi làm việc với quy trình thực hành an toàn và giải thích vì sao phải tuân thủ.
\loigiai{
Nguyên tắc cốt lõi là lập trình tự chuẩn bị, thao tác, giám sát, kết thúc và xử lí sự cố. Phải tuân thủ vì bảo đảm công việc có kiểm soát từ trước, trong và sau thí nghiệm; đồng thời luôn dừng và báo người phụ trách khi điều kiện không bảo đảm.
}
\end{bt}

% ===== Câu 121 =====
% ID: L10C1B2-37-TLN
% Mức: NB
\begin{ex}
% ID: L10C1B2-37-TLN
% Mức: NB
Một quy trình về quy trình thực hành an toàn gồm: (1) nhận diện nguy cơ; (2) lập trình tự chuẩn bị, thao tác, giám sát, kết thúc và xử lí sự cố; (3) chỉ ghi các bước thao tác mà không có bước kiểm tra và ứng phó; (4) báo giáo viên khi có bất thường; (5) chỉ tiếp tục khi nguy cơ đã được kiểm soát. Có bao nhiêu bước phù hợp quy tắc an toàn? Chỉ ghi số.
\shortans{4}
\loigiai{
Các bước (1), (2), (4), (5) phù hợp; bước (3) không an toàn. Có 4 bước phù hợp.
}
\end{ex}

% ===== Câu 125 =====
% ID: L10C1B2-39-TL
% Mức: VD
\begin{bt}
% ID: L10C1B2-39-TL
% Mức: VD
Xây dựng một bảng kiểm ngắn gồm ít nhất bốn mục cho tình huống khi giáo viên yêu cầu dừng thao tác liên quan đến quy trình thực hành an toàn.
\loigiai{
Bảng kiểm gồm: đọc hướng dẫn; nhận diện mối nguy; kiểm tra dụng cụ và bảo hộ; thực hiện lập trình tự chuẩn bị, thao tác, giám sát, kết thúc và xử lí sự cố; theo dõi bất thường; thu dọn và báo cáo sau thí nghiệm.
}
\end{bt}

% ===== Câu 128 =====
% ID: L10C1B2-40-TL
% Mức: VD
\begin{bt}
% ID: L10C1B2-40-TL
% Mức: VD
Một học sinh đã chỉ ghi các bước thao tác mà không có bước kiểm tra và ứng phó. Hãy nêu trình tự xử lí ngay sau khi phát hiện hành vi này.
\loigiai{
Trình tự: yêu cầu dừng ngay; không tiếp cận nếu chưa an toàn; cô lập nguồn nguy hiểm; báo giáo viên; sơ cứu hoặc xử lí theo hướng dẫn; chỉ hoạt động lại sau khi được xác nhận an toàn.
}
\end{bt}

% ===== Câu 131 =====
% ID: L10C1B2-42-TL
% Mức: VDC
\begin{bt}
% ID: L10C1B2-42-TL
% Mức: VDC
Đề xuất cách tổ chức nhóm học sinh để thực hiện nhiệm vụ về quy trình thực hành an toàn an toàn và có thể kiểm tra được.
\loigiai{
Phân công trưởng nhóm kiểm tra điều kiện, người thao tác, người quan sát an toàn và người ghi chép. Dùng bảng kiểm, xác nhận trước mỗi bước, quy định tín hiệu dừng và báo ngay bất thường.
}
\end{bt}

% ===== Câu 144 =====
% ID: L10C1B2-49-DS
% Mức: NB
\begin{ex}
% ID: L10C1B2-49-DS
% Mức: NB
Trong tình huống khi có người không đeo phương tiện bảo hộ đối với quy trình thực hành an toàn, hãy đánh giá các phát biểu sau.
\choiceTF
{\True Cần lập trình tự chuẩn bị, thao tác, giám sát, kết thúc và xử lí sự cố.}
{Có thể chỉ ghi các bước thao tác mà không có bước kiểm tra và ứng phó nếu thao tác diễn ra nhanh.}
{\True Phải báo cho giáo viên hoặc người phụ trách khi phát hiện nguy cơ vượt khả năng kiểm soát.}
{\True Chỉ tiếp tục thí nghiệm sau khi nguyên nhân nguy hiểm đã được loại bỏ hoặc kiểm soát.}
\loigiai{
\begin{itemchoice}
\item (Đúng) Cần lập trình tự chuẩn bị, thao tác, giám sát, kết thúc và xử lí sự cố. (Vì): vì đây là biện pháp an toàn của dạng quy trình thực hành an toàn; b sai vì hành vi “chỉ ghi các bước thao tác mà không có bước kiểm tra và ứng phó” vẫn nguy hiểm; c và d đúng do sự cố phải được báo cáo, cô lập và kiểm soát trước khi tiếp tục
\item (Sai) Có thể chỉ ghi các bước thao tác mà không có bước kiểm tra và ứng phó nếu thao tác diễn ra nhanh.
\item (Đúng) Phải báo cho giáo viên hoặc người phụ trách khi phát hiện nguy cơ vượt khả năng kiểm soát.
\item (Đúng) Chỉ tiếp tục thí nghiệm sau khi nguyên nhân nguy hiểm đã được loại bỏ hoặc kiểm soát.
\end{itemchoice}
}
\end{ex}

% ===== Câu 147 =====
% ID: L10C1B2-51-DS
% Mức: TH
\begin{ex}
% ID: L10C1B2-51-DS
% Mức: TH
Trong tình huống khi giáo viên yêu cầu dừng thao tác đối với quy trình thực hành an toàn, hãy đánh giá các phát biểu sau.
\choiceTF
{\True Cần lập trình tự chuẩn bị, thao tác, giám sát, kết thúc và xử lí sự cố.}
{Có thể chỉ ghi các bước thao tác mà không có bước kiểm tra và ứng phó nếu thao tác diễn ra nhanh.}
{\True Phải báo cho giáo viên hoặc người phụ trách khi phát hiện nguy cơ vượt khả năng kiểm soát.}
{\True Chỉ tiếp tục thí nghiệm sau khi nguyên nhân nguy hiểm đã được loại bỏ hoặc kiểm soát.}
\loigiai{
\begin{itemchoice}
\item (Đúng) Cần lập trình tự chuẩn bị, thao tác, giám sát, kết thúc và xử lí sự cố. (Vì): vì đây là biện pháp an toàn của dạng quy trình thực hành an toàn; b sai vì hành vi “chỉ ghi các bước thao tác mà không có bước kiểm tra và ứng phó” vẫn nguy hiểm; c và d đúng do sự cố phải được báo cáo, cô lập và kiểm soát trước khi tiếp tục
\item (Sai) Có thể chỉ ghi các bước thao tác mà không có bước kiểm tra và ứng phó nếu thao tác diễn ra nhanh.
\item (Đúng) Phải báo cho giáo viên hoặc người phụ trách khi phát hiện nguy cơ vượt khả năng kiểm soát.
\item (Đúng) Chỉ tiếp tục thí nghiệm sau khi nguyên nhân nguy hiểm đã được loại bỏ hoặc kiểm soát.
\end{itemchoice}
}
\end{ex}

% ===== Câu 149 =====
% ID: L10C1B2-53-DS
% Mức: VD
\begin{ex}
% ID: L10C1B2-53-DS
% Mức: VD
Trong tình huống khi phát hiện thiết bị có dấu hiệu bất thường đối với quy trình thực hành an toàn, hãy đánh giá các phát biểu sau.
\choiceTF
{\True Cần lập trình tự chuẩn bị, thao tác, giám sát, kết thúc và xử lí sự cố.}
{Có thể chỉ ghi các bước thao tác mà không có bước kiểm tra và ứng phó nếu thao tác diễn ra nhanh.}
{\True Phải báo cho giáo viên hoặc người phụ trách khi phát hiện nguy cơ vượt khả năng kiểm soát.}
{\True Chỉ tiếp tục thí nghiệm sau khi nguyên nhân nguy hiểm đã được loại bỏ hoặc kiểm soát.}
\loigiai{
\begin{itemchoice}
\item (Đúng) Cần lập trình tự chuẩn bị, thao tác, giám sát, kết thúc và xử lí sự cố. (Vì): vì đây là biện pháp an toàn của dạng quy trình thực hành an toàn; b sai vì hành vi “chỉ ghi các bước thao tác mà không có bước kiểm tra và ứng phó” vẫn nguy hiểm; c và d đúng do sự cố phải được báo cáo, cô lập và kiểm soát trước khi tiếp tục
\item (Sai) Có thể chỉ ghi các bước thao tác mà không có bước kiểm tra và ứng phó nếu thao tác diễn ra nhanh.
\item (Đúng) Phải báo cho giáo viên hoặc người phụ trách khi phát hiện nguy cơ vượt khả năng kiểm soát.
\item (Đúng) Chỉ tiếp tục thí nghiệm sau khi nguyên nhân nguy hiểm đã được loại bỏ hoặc kiểm soát.
\end{itemchoice}
}
\end{ex}

